\section{Evaluation}
\label{sec:evidence}
We characterize pre-barrier costs in Direct microbenchmarks and their contribution to MoE inference.

\subsection{Experimental Setup and Methodology}
\label{sec:setup}
Each node has eight H200 GPUs connected by NVSwitch and eight 400\,GbE RoCE interfaces. Cross-node runs use GIN's GDAKI backend, traffic class 162, PyTorch 2.13, CUDA 13, NCCL 2.31.2, and DeepEP V2 based on commit \texttt{01dc3aaa}~\cite{deepepcode}. Direct is selected with Hybrid disabled; available local communication within Direct is retained.

Microbenchmarks fix $H=7168$, $\topk=8$, 256 experts, BF16, and balanced uniform routing. $T$ denotes tokens per rank. Each point has three independent process runs with 200 active iterations per run. Error bars show the standard deviation across run means.

For each iteration, we select the rank with the shortest pre-barrier interval and divide that duration by the same rank's trace-on operator time. We average these ratios within each run and then across runs. This statistic reduces the effect of rank arrival skew, while retaining any waiting experienced by the selected rank. Phase shares use paired trace-on measurements; trace-off operator timings are kept separately.

The EP8 share sweep covers $T=8,32,64,128$. The EP-size comparison uses the original measurements at $T=8,32,96,128$, supplemented by the later EP8 $T=64$ batch. Configuration and per-run values accompany the data.

\subsection{Pre-Barrier Cost}
\label{sec:entry-eval}
Figure~\ref{fig:entry-cost} reports pre-barrier duration and operator share at EP8. Dispatch and Combine remain near 15\us{} across the measured token counts. Their shares fall from 20.57\% and 19.84\% at $T=8$ to 4.88\% and 5.05\% at $T=128$. Small batches therefore expose a larger fraction of the operation to readiness coordination.

% Frozen values: data/h200/numbers/paired_control_share.json.
\begin{figure*}[t]
\centering
\begin{minipage}{0.49\textwidth}
\centering
\begin{tikzpicture}
\begin{axis}[width=\linewidth,height=5.0cm,
 title={(a) Entry duration},xlabel={Tokens per rank},ylabel={Duration (\(\mu\mathrm{s}\))},
 xmode=log,log basis x=2,xtick={8,32,64,128},xticklabels={8,32,64,128},
 xmin=6,xmax=170,ymin=0,ymax=18,
 grid=major,grid style={black!10},axis lines=left,
 tick label style={font=\footnotesize},label style={font=\footnotesize},
 title style={font=\small},legend style={font=\footnotesize,at={(.97,.12)},anchor=south east,draw=none}]
\addplot+[datablue,mark=*,solid,error bars/.cd,y dir=both,y explicit]
 coordinates {(8,15.09) +- (0,0.06) (32,15.35) +- (0,0.46) (64,14.90) +- (0,0.10) (128,15.05) +- (0,0.21)};
\addlegendentry{Dispatch}
\addplot+[controlorange,mark=square*,dashed,error bars/.cd,y dir=both,y explicit]
 coordinates {(8,15.15) +- (0,0.03) (32,15.04) +- (0,0.22) (64,14.82) +- (0,0.13) (128,14.95) +- (0,0.23)};
\addlegendentry{Combine}
\end{axis}
\end{tikzpicture}
\end{minipage}
\hfill
\begin{minipage}{0.49\textwidth}
\centering
\begin{tikzpicture}
\begin{axis}[width=\linewidth,height=5.0cm,
 title={(b) Share of operator},xlabel={Tokens per rank},ylabel={Entry share (\%)},
 xmode=log,log basis x=2,xtick={8,32,64,128},xticklabels={8,32,64,128},
 xmin=6,xmax=170,ymin=0,ymax=24,
 grid=major,grid style={black!10},axis lines=left,
 tick label style={font=\footnotesize},label style={font=\footnotesize},
 title style={font=\small},legend style={font=\footnotesize,at={(.97,.97)},anchor=north east,draw=none}]
\addplot+[datablue,mark=*,solid,error bars/.cd,y dir=both,y explicit]
 coordinates {(8,20.57) +- (0,0.06) (32,12.25) +- (0,0.17) (64,8.29) +- (0,0.51) (128,4.88) +- (0,0.04)};
\addlegendentry{Dispatch}
\addplot+[controlorange,mark=square*,dashed,error bars/.cd,y dir=both,y explicit]
 coordinates {(8,19.84) +- (0,0.13) (32,12.26) +- (0,0.26) (64,8.20) +- (0,0.16) (128,5.05) +- (0,0.19)};
\addlegendentry{Combine}
\end{axis}
\end{tikzpicture}
\end{minipage}
\caption{H200 cross-node EP8 Direct. Entry duration stays close to 15\us{} while its operator share decreases. Both use the selected rank's trace-on measurements; error bars are standard deviations across three run means. Lines connect measured points.}
\label{fig:entry-cost}
\Description{Two panels show entry synchronization versus tokens per rank. Duration stays near fifteen microseconds for Dispatch and Combine, while operator share declines from about twenty percent to about five percent. Error bars describe three process runs.}
\end{figure*}


\subsection{Dependence on EP Size}
\label{sec:scale-eval}
Table~\ref{tab:entry-constants} summarizes pre-barrier duration with a separate constant for EP4, EP8, and EP16. All listed token counts participate in the estimate. The largest deviation of a per-token-count mean from its configuration's constant is 0.16\us{}, supporting a nearly constant cost over the measured batch range.

\begin{table}[t]
\centering
\small
\setlength{\tabcolsep}{3pt}
\begin{tabular}{@{}llrrr@{}}
\toprule
EP & Fitted $T$ & Dispatch & Combine & Max. dev.\\
\midrule
4 & 8/32/96/128 & 11.65 & 11.54 & 0.06\\
8 & 8/32/64/96/128 & 14.96 & 14.94 & 0.16\\
16 & 8/32/96/128 & 18.36 & 18.28 & 0.09\\
\bottomrule
\end{tabular}
\caption{Pre-barrier constants and maximum deviation of per-$T$ run means, in \(\mu\mathrm{s}\). Each $T$ has three process runs. EP8 includes a later $T=64$ batch with different instrumentation.}
\label{tab:entry-constants}
\end{table}

The fitted duration rises from approximately 11.6\us{} at EP4 to 18.3\us{} at EP16. These configurations change GPUs per node, traffic destinations, and network concurrency together, so the observed variation reflects their combined effect.

\subsection{Pre-Barriers in MoE Inference}
\label{sec:e2e}
We measure Qwen3-30B-A3B in BF16 with vLLM 0.27 and DeepEP V2 Direct. The model has $H=2048$, 128 experts, and $\topk=8$. Runs use actual routing, TP=1, EP8, and requests with 128 input and 128 output tokens. We compare one node with eight ranks to two nodes with four ranks each, at concurrency 64 and 128.

The serving build includes tracing and an expert-padding correctness fix. Each sampled active rank has 48 Dispatch and 48 Combine trace entries per decode step. We sum its pre-barrier intervals, divide by that rank's GPU-step time, and average over ranks and three client windows on the same server instance.

% Source: outline_shares.json, e2e[].pre_share_percent only.
\begin{figure}[t]
\centering
\begin{tikzpicture}
\begin{axis}[width=\columnwidth,height=5.4cm,
 ybar,bar width=12pt,symbolic x coords={64,128},xtick=data,
 xlabel={Request concurrency},ylabel={Pre-barrier share of GPU step (\%)},
 enlarge x limits=0.45,ymin=0,ymax=22,
 grid=major,grid style={black!10},
 tick label style={font=\footnotesize},label style={font=\footnotesize},
 legend style={font=\footnotesize,at={(.5,.98)},anchor=north,legend columns=2,draw=none}]
\addplot+[fill=datablue!45,draw=datablue,error bars/.cd,y dir=both,y explicit]
 coordinates {(64,9.33) +- (0,0.12) (128,7.94) +- (0,0.15)};
\addlegendentry{1 node}
\addplot+[pattern=north east lines,pattern color=controlorange,draw=controlorange,error bars/.cd,y dir=both,y explicit]
 coordinates {(64,16.72) +- (0,0.10) (128,15.19) +- (0,0.05)};
\addlegendentry{2 nodes}
\end{axis}
\end{tikzpicture}
\caption{Qwen3-30B-A3B decode at fixed EP8. Pre-barrier intervals for Dispatch and Combine are paired with the same rank's GPU step. Error bars describe three client windows, not server restarts.}
\label{fig:inference-share}
\Description{Grouped bars compare one-node and two-node EP8 decode at request concurrency sixty-four and one hundred twenty-eight. Single-node pre-barrier shares are around eight to nine percent, and cross-node shares around fifteen to seventeen percent.}
\end{figure}


At concurrency 64, the pre-barrier share is 9.33\% within one node and 16.72\% across two nodes. At concurrency 128, the corresponding values are 7.94\% and 15.19\%. Thus, readiness coordination remains visible in the full decode step with real routing and model computation.

These percentages describe sampled GPU steps. Unlike the microbenchmark statistic, they average over ranks rather than selecting the shortest interval. Together, the two experiments show both the operator-level cost and its place in inference.
